\section{Прямые и обратные маршруты расчёта}

Практическая применимость модели определяется набором маршрутов расчёта.
Маршрут задаётся входной парой параметров
и определяет последовательность вычислительных шагов.
В проекте используются как прямые зависимости IF97,
так и обратные постановки с аналитическими и численными процедурами.

\subsection{Прямой маршрут по параметрам \texorpdfstring{$(p, T)$ и $(\rho, T)$}{(p, T) и (rho, T)}}

Прямой маршрут $(p, T)$ является базовым для IF97.
Для него, а также для естественной для региона~3 постановки $(\rho, T)$,
расчёт строится вокруг прямого использования фундаментальных уравнений состояния.
По входной паре $(p, T)$ выполняются следующие этапы:
\begin{enumerate}
\item валидация диапазонов входных значений в пределах области применимости стандарта;
\item определение региона IF97 по топологическим границам (линия насыщения, граница B23 и др.);
\item вычисление безразмерных переменных и значений потенциала в выбранном регионе;
\item расчёт производных потенциала и итоговых свойств;
\item формирование унифицированного результата, включающего свойства и номер региона.
\end{enumerate}

В регионе~3 маршрут $(p, T)$ включает шаг определения плотности,
поскольку уравнения заданы в переменных $(\rho, T)$.
Для этого используются обратные зависимости $v(p, T)$,
опубликованные IAPWS в дополнительном релизе~\cite{iapws_vpt_region3},
и контроль физической допустимости найденного состояния.

Если же входная пара сразу задана как $(\rho, T)$,
то для региона~3 свойства вычисляются непосредственно
через потенциал Гельмгольца без дополнительной смены переменных.
В однофазных регионах 1, 2 и 5
маршрут $(\rho, T)$ трактуется как расширенная инженерная постановка:
по заданным плотности и температуре сначала восстанавливается давление,
после чего вычисление продолжается по стандартному маршруту $(p, T)$.

\subsection{Итерационные методы для маршрутов \texorpdfstring{$(p, h)$, $(p, s)$ и $(p, x)$}{(p, h), (p, s) и (p, x)}}

Обратные маршруты применяются,
когда состояние задаётся энергетическими параметрами.
К ним относятся входные пары $(p, h)$ и $(p, s)$.
В этих постановках нужно восстановить температуру
и далее вычислить полный набор свойств.

Для регионов и подрегионов,
где IF97 предоставляет аналитические обратные формулы,
температура определяется без итераций.
Это снижает вычислительные затраты
в массовых инженерных сценариях.

В регионе~1 используются единые зависимости
$T(p, h)$ и $T(p, s)$.
В регионе~2 стандарт задаёт отдельные семейства формул
для подобластей 2a, 2b и 2c.
Это позволяет учитывать сложную форму паровой области
без полного итерационного решения.
Для региона~5 аналитические обратные уравнения
в IF97 отсутствуют,
поэтому при попадании состояния в высокотемпературную область
требуется численное восстановление неизвестной температуры.

Если постановка $(p, h)$ или $(p, s)$ попадает в регион~3,
возникает система нелинейных уравнений относительно $(\rho, T)$.
В проекте для таких 2D-задач применяется решатель
Левенберга--Марквардта с предварительным сеточным поиском
и демпфированием шага.
Этот подход выбран для устойчивой сходимости
в околокритической области.

Маршрут $(p, x)$ используется для двухфазной области.
Хотя для него не требуется внешний итерационный контур,
он рассматривается вместе с обратными маршрутами,
поскольку также восстанавливает полное состояние
по неполному набору независимых параметров.
По заданным давлению $p$ и степени сухости $x$ определяется температура насыщения $T_{\mathrm{sat}}(p)$ и свойства насыщенной жидкости
($x=0$) и насыщенного пара ($x=1$).
В нижней части линии насыщения
эти свойства берутся из регионов~1 и~2,
а в околокритической части ---
из формулировки региона~3.
Для величин, аддитивных по массе, применяются правила смешения:
\begin{equation}
y = (1-x)\,y_{f} + x\,y_{g},
\end{equation}
\begin{gostwhere}
\item индекс $f$ соответствует насыщенной жидкости;
\item индекс $g$ --- насыщенному пару.
\end{gostwhere}

Удельный объём смеси вычисляется по этой формуле, а плотность определяется как $\rho = 1/v$.
Свойства двухфазной области, включая поведение $c_p$ и $w$,
требуют отдельной интерпретации в прикладных расчётах.
